% =============================================================================
%  Acelerador em FPGA do Supervisor de Seguranca LiDAR de uma Cadeira de Rodas
%  Ponto de Controle 1 (PC1) -- Proposta e Especificacao
%  PCR / Co-Projeto Hardware-Software -- UnB 2026/1
%
%  Compilacao: pdflatex relatorio_pcr_pc1.tex  (2x, por causa das referencias)
%  Requer figs/cobertura.dat (gerado por golden/gen_figs.py).
%
%  Este e o PRIMEIRO ponto de controle: define o problema, escolhe o bloco a
%  acelerar, especifica o algoritmo e propoe a arquitetura. Todos os numeros de
%  custo de hardware sao ESTIMATIVAS ANALITICAS -- ainda nao ha RTL, simulacao
%  nem sintese. Esses resultados chegam em PC2 (implementacao/verificacao) e
%  PC3 (sintese e fechamento).
% =============================================================================
\documentclass[conference]{IEEEtran}
\IEEEoverridecommandlockouts

\usepackage[utf8]{inputenc}
\usepackage[T1]{fontenc}
\usepackage[brazilian]{babel}
\usepackage{graphicx}
\usepackage{amsmath,amssymb}
\usepackage{booktabs}
\usepackage{cite}
\usepackage{url}
\usepackage{xcolor}
\usepackage{tikz}
\usepackage{pgfplots}
\usepackage{listings}
\usepackage{siunitx}
\pgfplotsset{compat=1.17}
\usetikzlibrary{arrows.meta, positioning, fit, calc}

\sisetup{output-decimal-marker={,}, per-mode=symbol, detect-all}

\lstset{basicstyle=\ttfamily\scriptsize, breaklines=true, columns=fullflexible,
        frame=single, framesep=3pt, xleftmargin=2pt}

\begin{document}

\title{Acelerador em FPGA para o Supervisor de Seguran\c{c}a\\
LiDAR de uma Cadeira de Rodas Semi-Assistida:\\
\large Ponto de Controle 1 --- Proposta, Especifica\c{c}\~ao e Particionamento}

\author{
\IEEEauthorblockN{Vítor Gonçalves de Andrade Silva}
\IEEEauthorblockA{Engenharia Eletrônica --- FCTE/UnB\\
Email: 261025182@aluno.unb.br}
\and
\IEEEauthorblockN{Thiago Henrique Oliveira Souza}
\IEEEauthorblockA{Engenharia Eletrônica --- FCTE/UnB\\
Email: 180131672@aluno.unb.br}
}

\maketitle

% =============================================================================
\begin{abstract}
Este primeiro ponto de controle define o escopo de um co-projeto
hardware--software sobre Zynq-7000 e justifica, por medi\c{c}\~ao, o bloco a ser
acelerado em FPGA. O sistema alvo \'e uma cadeira de rodas motorizada
semi-assistida cujo supervisor de seguran\c{c}a, rodando em Raspberry~Pi~5 com
ROS\,2, decide a cada varredura de LiDAR se o avan\c{c}o comandado pelo usu\'ario
deve ser atenuado. Perfilamos o supervisor e identificamos o teste de corredor
--- proje\c{c}\~ao dos \num{720} feixes e avalia\c{c}\~ao de \num{19} dire\c{c}\~oes
candidatas, \num{9918} avalia\c{c}\~oes por varredura --- como o \'unico bloco de
custo $O(N\cdot K)$ e o \'unico cuja lat\^encia tem consequ\^encia f\'isica
direta. Propomos um particionamento em que esse custo migra para hardware e a
fun\c{c}\~ao custo $O(K)$, onde residem os ganhos ajust\'aveis em campo,
permanece no ARM. Especificamos a arquitetura --- um CORDIC pipelinado
compartilhado alimentando \num{19} pistas paralelas de coeficiente constante ---,
os formatos de ponto fixo e as interfaces AXI4. Todos os n\'umeros de custo de
hardware apresentados aqui s\~ao \textbf{estimativas anal\'iticas}; a
implementa\c{c}\~ao RTL, a verifica\c{c}\~ao por equival\^encia e a s\'intese
s\~ao objeto dos pr\'oximos pontos de controle. Estabelecemos ainda o protocolo
de verifica\c{c}\~ao por modelo de ouro e o cronograma at\'e a entrega final.
\end{abstract}

% =============================================================================
\section{Introdu\c{c}\~ao e Motiva\c{c}\~ao}

Cadeiras de rodas motorizadas s\~ao inacess\'iveis a uma parcela expressiva dos
usu\'arios que mais precisariam delas: pesquisa cl\'inica cl\'assica aponta que
profissionais de reabilita\c{c}\~ao consideram muitos pacientes inaptos ao uso
justamente pela inadequa\c{c}\~ao das interfaces de controle~\cite{fehr2000}.
Espasmos e comandos acidentais no joystick produzem acelera\c{c}\~oes repentinas
e colis\~oes.

O projeto no qual este trabalho se insere converte uma cadeira comercial em um
ve\'iculo semi-assistido sob um princ\'ipio de seguran\c{c}a expl\'icito: o
sistema de alto n\'ivel pode \textbf{frear ou desviar}, mas \textbf{nunca
acelerar}; a iniciativa de movimento \'e sempre do usu\'ario. Um supervisor
executando em Raspberry~Pi~5 com ROS\,2~\cite{ros2docs} l\^e um LiDAR RPLIDAR~C1 e
atenua o comando do joystick antes de repass\'a-lo ao firmware de tempo real em
ESP32-S3. Este co-projeto n\~ao reescreve o sistema: escolhe o seu bloco mais
cr\'itico e o leva para hardware.

\subsection{O bloco a acelerar e por qu\^e}

A escolha do bloco n\~ao \'e arbitr\'aria --- \'e resultado de perfilamento do
supervisor j\'a existente em software. O teste de corredor \'e, disparadamente, o
n\'o mais pesado: medimos \SI{17.7}{\percent} da CPU da Pi~5 na implementa\c{c}\~ao
escalar e \SI{8.3}{\percent} ap\'os vetoriza\c{c}\~ao com NumPy. Mais importante
que o custo, por\'em, \'e a natureza da sua lat\^encia: \emph{ela tem
consequ\^encia f\'isica}. O supervisor s\'o corta o comando de avan\c{c}o; a
cadeira ainda percorre uma dist\^ancia por in\'ercia. Todo atraso entre a chegada
da varredura e a decis\~ao se traduz em cent\'imetros adicionais rumo ao
obst\'aculo. Trata-se, portanto, de um caso em que acelerar n\~ao \'e apenas
otimiza\c{c}\~ao, e sim seguran\c{c}a --- o que motiva o co-projeto.

A escolha \'e ainda mais pertinente porque uma revis\~ao recente do supervisor
revelou uma falha geom\'etrica de seguran\c{c}a. A folga frontal era medida por um
\emph{cone angular} ancorado no sensor, constru\c{c}\~ao v\'alida apenas se o
LiDAR estivesse sobre o eixo longitudinal do ve\'iculo. Estando ele
\SI{0.336}{\metre} fora do eixo, o cone varre uma faixa deslocada e a metade
direita da cadeira fica cega justamente nas curtas dist\^ancias
(Fig.~\ref{fig:cobertura}). A corre\c{c}\~ao substitui o cone pelo teste do
\emph{corredor f\'isico} que a cadeira ocupar\'a --- e \'e esse algoritmo,
correto por\'em mais caro, que propomos acelerar.

\begin{figure}[!t]
\centering
\begin{tikzpicture}
\begin{axis}[width=\columnwidth, height=4.2cm,
  xlabel={dist\^ancia do obst\'aculo (m)},
  ylabel={cobertura (\%)}, ymin=0, ymax=105, xmin=0.3, xmax=1.5,
  legend pos=south east, legend style={font=\scriptsize},
  grid=major, grid style={gray!20}, tick label style={font=\scriptsize},
  label style={font=\scriptsize}]
  \addplot[thick, red!70!black] table[x=d, y=cone30] {figs/cobertura.dat};
  \addlegendentry{cone $\pm30^\circ$}
  \addplot[thick, orange, dashed] table[x=d, y=cone20] {figs/cobertura.dat};
  \addlegendentry{cone $\pm20^\circ$}
  \addplot[thick, blue!60!black, densely dotted] coordinates {(0.3,100) (1.5,100)};
  \addlegendentry{corredor (proposto)}
\end{axis}
\end{tikzpicture}
\caption{Fra\c{c}\~ao da largura da cadeira efetivamente monitorada. O cone
ancorado no sensor \emph{piora} quanto mais perto est\'a o obst\'aculo, porque
estreita enquanto o desvio lateral do v\'ertice permanece constante: a
\SI{0.60}{\metre} cobre \SI{52}{\percent}, e a parcela n\~ao coberta \'e sempre a
metade direita. O teste de corredor cobre \SI{100}{\percent} por
constru\c{c}\~ao. \'E esta falha que motiva o algoritmo mais caro.}
\label{fig:cobertura}
\end{figure}

\subsection{Trabalhos relacionados}

A fun\c{c}\~ao custo que pontua dire\c{c}\~oes candidatas pela folga \'e
aparentada ao \emph{Vector Field Histogram}~\cite{borenstein1991}, com a
distin\c{c}\~ao essencial de raciocinar sobre a \emph{planta} do rob\^o e n\~ao
sobre um cone do sensor. O CORDIC~\cite{volder1959,andraka1998} \'e a escolha
cl\'assica para trigonometria em FPGA sem multiplicadores. O uso de aceleradores
acoplados a processadores embarcados via AXI4 sobre Zynq segue a pr\'atica
consolidada em~\cite{zynqbook}.

% =============================================================================
\section{Especifica\c{c}\~ao do Algoritmo}

Cada retorno do LiDAR, um par $(r_i, a_i)$, \'e projetado no referencial
\texttt{base\_link} (origem no centro do eixo traseiro):
\begin{align}
X_i &= x_L + r_i\cos(a_i + \theta_L), &
Y_i &= y_L + r_i\sin(a_i + \theta_L),
\label{eq:proj}
\end{align}
\noindent onde $(x_L, y_L, \theta_L)$ \'e a pose calibrada do sensor. Para avaliar
um candidato de desvio $\delta_k$, giram-se os pontos por $-\delta_k$ em torno do
eixo traseiro --- o centro instant\^aneo de rota\c{c}\~ao do acionamento
diferencial:
\begin{align}
x^{(k)}_i &= X_i\cos\delta_k + Y_i\sin\delta_k, \\
y^{(k)}_i &= -X_i\sin\delta_k + Y_i\cos\delta_k .
\label{eq:rot}
\end{align}
O ponto \'e obst\'aculo se $|y^{(k)}_i| \le W/2$ e $x^{(k)}_i > X_f$, com folga
$x^{(k)}_i - X_f$; a folga do candidato \'e o m\'inimo sobre os pontos. Aqui
$W=\SI{0.61}{\metre}$ \'e a largura da cadeira e $X_f=\SI{0.667}{\metre}$ a frente
do chassi. Note que a estrutura da pr\'opria cadeira cai \emph{atr\'as} de $X_f$ e
\'e descartada por constru\c{c}\~ao, sem nenhum limiar arbitr\'ario de
dist\^ancia m\'inima.

Com $N=\num{720}$ feixes e $K=\num{19}$ candidatos ($\delta_k$ de $-45^\circ$ a
$+45^\circ$, passo $5^\circ$), s\~ao \num{9918} avalia\c{c}\~oes por varredura ---
o custo $O(N\cdot K)$ que motiva o acelerador.

% =============================================================================
\section{Particionamento Hardware--Software Proposto}

O supervisor decomp\~oe-se em duas partes de custo muito distinto
(Tab.~\ref{tab:part}). Propomos migrar apenas a primeira para hardware.

\begin{table}[!t]
\caption{Particionamento HW/SW proposto}
\label{tab:part}
\centering\footnotesize
\begin{tabular}{@{}lll@{}}
\toprule
\textbf{Etapa} & \textbf{Custo} & \textbf{Onde} \\
\midrule
Proje\c{c}\~ao + teste de corredor & $O(N\!\cdot\!K)=\num{9918}$ & \textbf{FPGA} \\
Fun\c{c}\~ao custo + \emph{argmin} & $O(K)=\num{19}$ & ARM \\
\bottomrule
\end{tabular}
\end{table}

A decis\~ao de \emph{n\~ao} levar tudo para hardware \'e deliberada. A fun\c{c}\~ao
custo s\~ao \num{19} compara\c{c}\~oes: o ganho seria irris\'orio. Em compensa\c{c}\~ao,
\'e nela que residem os pesos ajustados em campo ($w_{\text{obs}}$,
$w_{\text{dev}}$, $K_{\text{assist}}$), retocados a cada ensaio. Congel\'a-los em
RTL trocaria flexibilidade por nada --- exatamente o \emph{trade-off} que o
co-projeto existe para arbitrar.

% =============================================================================
\section{Arquitetura Proposta}

A Fig.~\ref{fig:arq} mostra o \emph{datapath} que pretendemos implementar. O
CORDIC \'e \textbf{compartilhado} porque a proje\c{c}\~ao~\eqref{eq:proj} \'e
comum a todos os candidatos: o custo $O(N)$ \'e pago uma \'unica vez. O que se
replica s\~ao as \num{19} pistas, que carregam o custo $O(N\cdot K)$ --- e \'e
precisamente a\'i que esperamos que o paralelismo espacial da FPGA renda: as
pistas avaliam \emph{no mesmo ciclo} aquilo que o processador precisa iterar.

\begin{figure}[!t]
\centering
\begin{tikzpicture}[font=\scriptsize, node distance=3.5mm,
  blk/.style={draw, rounded corners=2pt, minimum height=7mm, inner sep=3pt,
              fill=gray!8, align=center},
  arr/.style={-Latex}]
  \node[blk] (in) {AXI4-Stream\\$(r_i,\theta_i)$};
  \node[blk, right=of in, fill=blue!8] (cor) {CORDIC\\16 est\'agios\\\textbf{0 DSP}};
  \node[blk, right=of cor] (prj) {proje\c{c}\~ao\\$X_i, Y_i$};
  \node[blk, right=6mm of prj, fill=orange!10, minimum height=13mm]
       (lanes) {19 pistas\\paralelas\\\textbf{$\sim$76 DSP}};
  \node[blk, below=5mm of lanes] (out) {AXI4-Lite\\19 folgas};
  \draw[arr] (in) -- (cor);
  \draw[arr] (cor) -- node[above, font=\tiny] {$\cos,\sin$} (prj);
  \draw[arr] (prj) -- (lanes);
  \draw[arr] (lanes) -- (out);
\end{tikzpicture}
\caption{\emph{Datapath} proposto. Vaz\~ao alvo de \textbf{1 feixe/ciclo}, sem
bolhas. O CORDIC \'e compartilhado (custo $O(N)$); as pistas s\~ao replicadas
(custo $O(N\cdot K)$). Os n\'umeros de DSP s\~ao \emph{estimativas} a confirmar na
s\'intese (PC3).}
\label{fig:arq}
\end{figure}

\paragraph{CORDIC}
Optamos por CORDIC em modo rota\c{c}\~ao~\cite{volder1959}, pipelinado em
\num{16} est\'agios, em vez de uma LUT de seno. Uma LUT exigiria uma ROM indexada
pelo \'indice do feixe, o que s\'o funciona se \texttt{angle\_min}, o incremento e
o \emph{yaw} forem fixos em s\'intese --- mas eles s\~ao \textbf{par\^ametros de
calibra\c{c}\~ao}, reescritos em tempo de execu\c{c}\~ao por AXI4-Lite. O CORDIC
aceita qualquer \^angulo e, por usar apenas somadores e deslocamentos cabeados,
n\~ao deve consumir \emph{nenhum} DSP48. Como o la\c{c}o s\'o converge para
$|z|\le\SI{1.7434}{\radian}$, um est\'agio de redu\c{c}\~ao de quadrante trar\'a
$\theta$ para $[-\pi/2,\pi/2]$ e inverter\'a o sinal na sa\'ida.

\paragraph{Pistas}
Em cada pista, $\cos\delta_k$ e $\sin\delta_k$ ser\~ao \textbf{constantes de
s\'intese} (\emph{generics}): a rota\c{c}\~ao~\eqref{eq:rot} vira quatro
multiplica\c{c}\~oes de coeficiente fixo, mape\'aveis diretamente em DSP48.

\subsection{Formato de ponto fixo}

Os formatos propostos est\~ao na Tab.~\ref{tab:fmt}. A escolha da \emph{faixa} ---
e n\~ao apenas da resolu\c{c}\~ao --- \'e um ponto de aten\c{c}\~ao que
pretendemos validar cuidadosamente na verifica\c{c}\~ao de PC2.

\begin{table}[!t]
\caption{Formatos de ponto fixo propostos}
\label{tab:fmt}
\centering\footnotesize
\begin{tabular}{@{}llll@{}}
\toprule
\textbf{Grandeza} & \textbf{Formato} & \textbf{Faixa} & \textbf{LSB} \\
\midrule
dist\^ancias, coords. & Q6.10 & $\pm\num{32}$\,m & \SI{0.98}{\milli\metre} \\
$\cos$, $\sin$ & Q1.14 & $\pm 2$ & \num{6.1e-5} \\
\^angulos & Q3.13 & $\pm\SI{4.0}{\radian}$ & \SI{1.22e-4}{\radian} \\
\bottomrule
\end{tabular}
\end{table}

\subsection{Interfaces AXI4}

Os feixes chegar\~ao por \textbf{AXI4-Stream} (\texttt{TDATA} $=$ [$\theta$(31:16)
$|$ $r$(15:0)], \texttt{TUSER} $=$ feixe v\'alido, \texttt{TLAST} $=$ fim da
varredura), alimentado por um AXI-DMA que l\^e a varredura da DDR. O controle e o
resultado usar\~ao \textbf{AXI4-Lite}: \texttt{CTRL} (\emph{start}),
\texttt{STATUS} (\emph{done}/\emph{busy}) e \num{19} registradores de folga. Uma
interrup\c{c}\~ao sinalizar\'a a conclus\~ao, dispensando espera ocupada no ARM.
A escolha \'e natural: \emph{stream} para o fluxo volumoso e sequencial (a
varredura), mapeado em mem\'oria para o punhado de escalares (calibra\c{c}\~ao e
resultado).

% =============================================================================
\section{Estimativas Anal\'iticas de Custo}

Como ainda n\~ao h\'a RTL nem s\'intese, os n\'umeros abaixo s\~ao previstos por
contagem de operadores e profundidade de \emph{pipeline}; ser\~ao confrontados
com a medi\c{c}\~ao em PC2 (lat\^encia por simula\c{c}\~ao) e PC3 (recursos e
\emph{timing} por s\'intese).

\begin{table}[!t]
\caption{Estimativas anal\'iticas (a confirmar em PC2/PC3)}
\label{tab:est}
\centering\footnotesize
\begin{tabular}{@{}lll@{}}
\toprule
\textbf{Grandeza} & \textbf{Estimativa} & \textbf{Confirmado em} \\
\midrule
DSP48 (pistas + proje\c{c}\~ao) & $\sim$\num{78} ($19\times4+2$) & PC3 (s\'intese) \\
DSP48 (CORDIC)      & \num{0}                    & PC3 (s\'intese) \\
BRAM                & \num{0} (\emph{streaming}) & PC3 (s\'intese) \\
Vaz\~ao             & \num{1} feixe/ciclo        & PC2 (simula\c{c}\~ao) \\
Lat\^encia          & $\sim N +$ profundidade    & PC2 (simula\c{c}\~ao) \\
$F_{\max}$ alvo     & \SI{100}{\mega\hertz}      & PC3 (s\'intese) \\
\bottomrule
\end{tabular}
\end{table}

A lat\^encia prevista \'e da ordem de $N$ ciclos (um feixe por ciclo) mais a
profundidade do \emph{pipeline} do CORDIC e das pistas, algo em torno de
\num{750} ciclos por varredura; a \SI{100}{\mega\hertz}, aproximadamente
\SI{7.5}{\micro\second}. Contra o \emph{baseline} de software a medir, isso deve
render duas a tr\^es ordens de grandeza --- hip\'otese que PC2 verifica.

% =============================================================================
\section{Protocolo de Verifica\c{c}\~ao e Cronograma}

\subsection{Verifica\c{c}\~ao por modelo de ouro}
Estabelecemos desde j\'a a estrat\'egia de verifica\c{c}\~ao, para que a
implementa\c{c}\~ao nasça test\'avel. N\~ao verificaremos por inspe\c{c}\~ao de
formas de onda, e sim por \textbf{equival\^encia bit a bit} contra um modelo de
refer\^encia:
\begin{enumerate}
\item um \textbf{modelo em ponto fixo} (Python) que emula o RTL exatamente
(mesmo CORDIC, mesmos deslocamentos, mesmos truncamentos);
\item um \textbf{modelo em ponto flutuante} como refer\^encia f\'isica, para medir
o erro de quantiza\c{c}\~ao;
\item as constantes do RTL (\texttt{cordic\_pkg.vhd}) ser\~ao \textbf{geradas}
pelo pr\'oprio modelo, de modo que RTL e refer\^encia casem por constru\c{c}\~ao;
\item o est\'imulo n\~ao ser\'a sint\'etico: usaremos varreduras \textbf{reais} da
cadeira, capturadas com um obst\'aculo no antigo ponto cego.
\end{enumerate}

\subsection{Cronograma}
\begin{table}[!t]
\caption{Cronograma at\'e a entrega final}
\label{tab:crono}
\centering\footnotesize
\begin{tabular}{@{}lll@{}}
\toprule
\textbf{Ponto} & \textbf{Entrega} & \textbf{Estado} \\
\midrule
PC1 & Proposta, algoritmo, particionamento, arquitetura & \textbf{este documento} \\
PC2 & RTL, modelo de ouro, verifica\c{c}\~ao, lat\^encia & planejado \\
PC3 & S\'intese, \emph{timing}, pot\^encia, \emph{block design} & planejado \\
\bottomrule
\end{tabular}
\end{table}

% =============================================================================
\section{Conclus\~ao Parcial}

Este ponto de controle fecha a fase de concep\c{c}\~ao. Justificamos por
perfilamento a escolha do teste de corredor como bloco a acelerar, mostramos que
sua lat\^encia tem consequ\^encia f\'isica de seguran\c{c}a, especificamos o
algoritmo e propusemos um particionamento HW/SW, uma arquitetura CORDIC $+$
pistas e os formatos de ponto fixo e interfaces AXI4. Deliberadamente n\~ao
apresentamos resultados medidos: as estimativas de recursos, vaz\~ao e
lat\^encia s\~ao anal\'iticas e ser\~ao confrontadas com implementa\c{c}\~ao e
s\'intese nos pr\'oximos pontos de controle. O risco t\'ecnico que mais nos
preocupa --- e que a verifica\c{c}\~ao de PC2 deve expor --- \'e a escolha da
faixa dos formatos de ponto fixo, sobretudo a dos \^angulos.

% =============================================================================
\begin{thebibliography}{9}

\bibitem{fehr2000}
L.~Fehr, W.~E. Langbein e S.~B. Skaar, ``Adequacy of power wheelchair control
interfaces for persons with severe disabilities: A clinical survey,''
\emph{J. Rehabil. Res. Dev.}, vol.~37, no.~3, pp.~353--360, 2000.

\bibitem{borenstein1991}
J.~Borenstein e Y.~Koren, ``The Vector Field Histogram --- Fast Obstacle
Avoidance for Mobile Robots,'' \emph{IEEE Trans. Robot. Autom.}, vol.~7, no.~3,
pp.~278--288, 1991.

\bibitem{volder1959}
J.~E. Volder, ``The CORDIC Trigonometric Computing Technique,'' \emph{IRE Trans.
Electron. Comput.}, vol.~EC-8, no.~3, pp.~330--334, 1959.

\bibitem{andraka1998}
R.~Andraka, ``A survey of CORDIC algorithms for FPGA based computers,'' em
\emph{Proc. ACM/SIGDA Int. Symp. FPGAs}, 1998, pp.~191--200.

\bibitem{zynqbook}
L.~H. Crockett, R.~A. Elliot, M.~A. Enderwitz e R.~W. Stewart, \emph{The Zynq
Book: Embedded Processing with the ARM Cortex-A9 on the Xilinx Zynq-7000}.
Glasgow: Strathclyde Academic Media, 2014.

\bibitem{ros2docs}
Open Robotics, ``ROS\,2 Jazzy Jalisco Documentation,'' 2024. [Online].
Dispon\'ivel: \url{https://docs.ros.org/en/jazzy/}

\end{thebibliography}

\end{document}
